\subsection{One coefficient space for the seed and the rational point}
\label{subsec:rh-four-common-space}

We supply the universal-space interface used in the propagation
argument, keeping the original coefficient and flag equations.
A complete ring of lifts of the diagonal residual representation
alone does not have the required universal property at the
nonsplit residual representations. The distinction between its
completed homogeneous pieces and the finite-type entry algebra is
essential.

\begin{theorem}[The original finite-type space in the RH towers]
\label{thm:rh-four-space}
Fix one of the original fitted deep RH towers, with frame
$\Phi=T.\mathrm{fr}(\varnothing)$, and a permitted complete
coefficient DVR $\mathcal O$ with finite residue field $k$.
There is a space $(B,A,\mathcal P)$ satisfying every field of
the original \texttt{Space}: in particular $B$ is complete
Noetherian local with residue $k$, $B\hookrightarrow A$,
$A$ is finite type over $B$, and the representation and flags
in $A$ induce every continuous original lift of every
\texttt{IsRes} pair over $k$. Moreover each original polynomial
residual link over $k[t]$ is induced by a map $A\to k[t]$
whose restriction to $B$ is the residue map.
\end{theorem}

\begin{proof}
Start with the complete framed universal ring $R$ of lifts
of the diagonal residual representation on $\Phi$,
imposing the prescribed symplectic similitude, the killing
set $N$, and diagonalization at the mark. Adjoin the complete
universal deformation rings of the independently labeled
local characters $\theta_{vi}$, initially without imposing
the flag equations or the local torsion conditions.

For completeness this ring exists and is Noetherian by the
following standard finite-generator construction. After
killing $N$, the relevant group is a quotient of the Galois
group unramified outside a finite set containing $\ell$:
outside the original ramification set every inertia element
belongs to $N$. The kernel of its finite diagonal residual
representation has finitely generated maximal pro-$\ell$
quotient, by finiteness of its first cohomology with
$\mathbb F_\ell$ coefficients (it is a quotient of the
corresponding finite-extension number-field Galois group,
with finitely many allowed ramified places). Every lift
to an Artinian local coefficient ring with finite residue
field factors through this quotient of the residual kernel.
The kernel quotient is characteristic and hence normal
in the full group. The resulting profinite group is
topologically finitely generated: combine generators of
the pro-$\ell$ kernel with lifts of generators of the
finite quotient. Choose these finitely many generators,
use formal matrix coordinates reducing to their given
residual values, and impose their profinite group relators.
In a finite quotient of the coefficient ring the matrices
are a finite group, so these relators have well-defined
compatible evaluations, giving equations in the complete
power series ring. The same construction for the local
characters uses local class field theory, or finiteness of
local first cohomology. Completion and passage to compatible
Artinian quotients represent lifts over arbitrary complete
Noetherian local rings with residue $k$. The closed
symplectic, killing, and marking equations are the literal
matrix-entry equations. This constructs $R$ and the
universal representation $D$ and independent characters
$\theta_{vi}$.

Give every representation entry weight
$\operatorname{wt}(i,j)$, and the local characters weight
zero. We can present $R$ by finitely many formal coordinates
of these weights: choose a finite subset of actual entry
differences from the residual diagonal lifts and of actual
character differences from their prescribed lifts,
generating the maximal ideal modulo its square and the
coefficient uniformizer. Such a subset exists because
the universal matrix and character coordinates topologically
generate $R$; Nakayama and completeness give the desired
surjective power series presentation. The entries on the
diagonal and all character differences have weight zero.
The constants used in the off-diagonal differences are
zero. The torus conjugation
$\operatorname{diag}(z_1z_2,z_2,z_1,1)$ acts on this
presentation, with $\mathcal O$ of weight zero, and
preserves the defining equations by the universal property.

Here and below a formal homogeneous piece is understood in
the completed sense, not as a polynomial. To justify the
grading in this topology, use $z_1=1+U,z_2=1+V$ in the
universal complete ring with two extra formal variables.
A relation in the presentation, after substitution, is zero
there. Modulo any power of the maximal ideal of $R$, only
finitely many monomials survive. Laurent monomials of
distinct bidegrees are linearly independent after substitution
of $1+U,1+V$ into a power series ring over any coefficient
ring (multiply by a common Laurent denominator and use
the invertible polynomial change of variables). Thus each
homogeneous part of the relation vanishes modulo every
power of the maximal ideal, and hence vanishes in $R$.
This yields the closed homogeneous relation ideal. Every
actual entry and local character is homogeneous by the same
argument applied to its transformation law. Write $R_w$ for
the image of the formal homogeneous series of weight $w$,
and set
\[
                 B_0=R_0,\qquad A_0=\bigoplus_{w}R_w
                    \ \subset R
\]
where the direct sum means finite sums, so $A_0$ is an
ordinary algebra rather than a completion. Independence
of the homogeneous pieces follows from the torus
substitution just used.

We recall the elementary monomial facts ensuring the
Noetherian and continuity properties of this construction.
The nonnegative exponent vectors of weight zero form a
finitely generated affine monoid (equivalently, take the
integer points of the rational polyhedral cone defined by
the weight-zero equations). For each fixed weight $w$,
Dickson's lemma gives finitely many minimal nonnegative
exponent vectors of that weight in the coordinatewise
order. Every other exponent vector of that weight is
one of these plus a weight-zero vector. Hence a homogeneous
formal series of weight $w$ is a sum of finitely many
corresponding monomials with coefficients formal series in
the finitely many generators of weight zero. This assertion
holds with the coefficient uniformizer included in the
topology: choose one such decomposition per exponent
vector; as total monomial degree tends to infinity, the
degree in weight-zero generators tends to infinity, since
their degrees and the degrees of the minimal representatives
are bounded. The same observation shows that for any
fixed $w$, the intersection of a sufficiently high power
of the presentation's maximal ideal with that homogeneous
piece maps into any prescribed power of the weight-zero
maximal ideal times $R_w$. In passing to the quotient one
may project a lift to its weight-$w$ component, because the
relation ideal is closed and homogeneous. These arguments
prove that $B_0$ is a quotient of a complete power series
ring on finitely many weight-zero monomials, local with
residue $k$, that the entry algebra $A_0$ is generated by
the finitely many original coordinates over $B_0$, and
that the matrix entries and the coefficient characters
are continuous in the $B_0$-adic topology in their
respective homogeneous pieces.

We next prove the required coefficient-generation
assertion over the actual integral ring, without dividing
by any factorial. Let $\lambda_i=D(g_0)_{ii}$. These four
elements of $B_0$ are determined by the characteristic
polynomial of the diagonal mark and their distinct labeled
residues, by Hensel's lemma. In particular their differences
are units in $B_0$. Products of matrix entries along
directed cycles are obtained as traces of products of
matrices and the diagonal spectral projectors of the mark.
Each projector is the explicit polynomial
\[
 E_i=\prod_{j\ne i}
       \frac{D(g_0)-\lambda_j}{\lambda_i-\lambda_j}.
\]
Expand the trace using this formula and the representation
law: it is a finite combination of traces of single group
elements with coefficients depending on the marked
eigenvalues and unit differences. The trace itself is the
negative next-to-leading characteristic-polynomial
coefficient.

Every weight-zero entry monomial reduces to such cycle
monomials. To see this directly in the rank-four root
system, draw an entry $D(g)_{ij}$ as an edge from $j$ to
$i$, counting repeated factors as separate edges.
The difference between the incoming and outgoing
multiplicities at the four vertices is a vector of the form
$(c,-c,-c,c)$, since its total and its two weight coordinates
are zero. Decompose the edges into directed cycles and
paths from vertices with excess outgoing multiplicity to
vertices with excess incoming multiplicity. If $c>0$,
a path goes from $\{1,2\}$ to $\{0,3\}$. Reverse the path
and replace each vertex by its complementary index
$i\mapsto3-i$. Every edge can be so replaced in the monomial:
the symplectic inverse identity identifies $D(g)_{ij}$ with
the entry of $D(g^{-1})$ at $(3-j,3-i)$ times a known unit
(the similitude and a sign). This changes the imbalance
vector by $(-1,1,1,-1)$, reducing $c$ by one. When $c<0$,
the same operation on a path from $\{0,3\}$ to $\{1,2\}$
increases $c$ by one. Iteration gives an ordinary balanced
directed graph, a disjoint union of cycles with edge
multiplicities. Thus the original monomial is a known
unit times a product of cycle traces. The argument is
an identity of matrix entries over the original coefficient
algebra, not just over a field of generic matrices.

Apply this result also to the coordinate differences
from the prescribed constant diagonal and character
lifts. Since their constants are of weight zero, their
balanced products are generated, with completion, by
the characteristic-polynomial coefficients of all group
elements and the original independent local characters.
Two local maps out of $B_0$ into the same complete local
coefficient ring agreeing on the latter data agree on
the four marked eigenvalues by their labeled residues,
on the cycle traces, on every balanced coordinate
monomial, and therefore on $B_0$ by continuity. This is
the exact \texttt{genB} uniqueness property.

Adjoin now to $A_0$ the entries of a lower symplectic
unitriangular flag for each original local place.
Impose the literal flag equations: the conjugated
local representations are upper triangular, their ordered
diagonal entries are the corresponding $\theta_{vi}$,
and the latter are trivial on $\mathrm{Mu}$. Take the
quotient $A$ by this ideal and let $B$ be the image of
$B_0$ therein. The ideal defining $B$ in $B_0$ is closed,
since $B_0$ is complete Noetherian local, so $B$ has
the same completeness property. The diagonal residual
pair with identity flags is one of the residue
evaluations and shows that this quotient is nonzero
and has the specified residue field. We now have
$B\hookrightarrow A$, finite type, and all matrix,
flag, continuity, characteristic-polynomial, and
coefficient uniqueness equations in \texttt{Space}.
In particular every character can be recovered from
the flag and representation entries, so those entries
generate $A$ over $B$.

The crucial remaining point is its lift universality.
Let $(r,n)$ be any original continuous lift of an
\texttt{IsRes} pair over $k$ to a complete local
coefficient ring $R'$ with finite residue $k$.
Take a linear functional $\tau$ on the root lattice
strictly positive on its allowed residual support:
$\tau=(-2,-1)$ for $nP$, and $\tau=(-1,-2)$
for $nM$. Its existence follows directly from the
four displayed roots in each support set.
Rescale the entries of $r$ by a formal cocharacter
$z^{\tau(w)}$, leaving the labeled local characters
fixed. To implement the rescaling integrally, for
every $m\ge1$ put $R_m'=R'/\mathfrak m_{R'}^m$ and
let $I_m$ be its nilpotent maximal ideal. In
$R_m'((z))$ take the subring
\[
      S_m=R_m'[[z]][\, I_m z^{-D}\,],
      \qquad D\ge \max_w|\tau(w)|.
\]
The adjoined generators are nilpotent, so this
is a finite module over $R_m'[[z]]$, complete local
with residue $k$ and maximal ideal generated by
$z$ and the adjoined nilpotent ideal together with
$I_m$. Every scaled matrix entry belongs to it:
an entry with negative weight was in $I_m$,
and each positive-weight entry is divisible by $z$.
Modulo its maximal ideal the rescaled representation
is the original split diagonal. Its multiplier,
killing and marking are unchanged; continuity holds
because the original entries and characters modulo
$\mathfrak m_{R'}^m$ are locally constant.
Thus the universal complete ring $R$ maps to $S_m$.

For each homogeneous series of weight $w$, its
image is a single Laurent monomial $z^{\tau(w)}$
times a coefficient in $R_m'$. Indeed this holds
on every monomial of a homogeneous lift, including
weight-zero constants, and their partial sums
converge by continuity of the universal local map.
The coefficient group $R_m' z^{\tau(w)}$ is a finite
closed subset of the separated ring $S_m$.
Consequently evaluation at $z=1$ is well-defined
on the image of the finite sums $A_0$, as those
images are Laurent polynomials. These homomorphisms
to $R_m'$ are compatible in $m$ by uniqueness of
the universal maps (the rings $S_m$ also map
compatibly under reduction). Their limit gives a
map $A_0\to R'$ sending the entries and independent
characters to the original values. Sending the
additional flag variables to the original flag
entries respects every imposed equation, yielding
$A\to R'$. Its restriction to $B$ is local: balanced
nonconstant monomials in the off-diagonal entries
cannot all have weights strictly positive for
$\tau$; at least one factor is in $\mathfrak m_{R'}$.
The diagonal and character coordinate differences
also belong to that ideal. Equivalently this follows
by taking the coefficient of $z^0$ in the above
construction modulo $\mathfrak m_{R'}$.
This proves universality for the original lift
functor. For a general complete Noetherian local
target the same proof applies since every Artinian
quotient with finite residue field is finite.

Finally take an actual polynomial residual pair
in the source \texttt{Linked} relation over $k[t]$.
Embed $k[t]$ into $k[[t]]$ and apply lift universality
to its constant term residual pair. The characteristic
polynomials throughout the link are the original
diagonal polynomials: each permitted support $nP$
or $nM$ is triangular for a fixed coordinate ordering.
The ordered local characters are also the original
residual characters, by the link condition.
The coefficient uniqueness just proved implies
that the induced map on $B$ is the constant residue
map. Since $A$ is algebraically generated over $B$
by the representation and flag entries, all of which
are polynomials in this map, its image lies in $k[t]$.
This is the required original line map, including
its constant coefficient map, not a path through
unrelated deformation spaces.
\end{proof}

\begin{corollary}[Simultaneous coefficient realization]
\label{cor:rh-four-common-coefficients}
In the preceding theorem one can choose $\mathcal O$ once
for the rational target, the finite residual chain, the
two ordinary links and the arithmetic seed, and realize
all of them in this one space. More precisely the rational
target is allowed to be of characteristic zero or
characteristic $\ell$ as in \texttt{Amb.Small}.
\end{corollary}
\begin{proof}
The original ambient residue field is algebraic over a
finite field by \texttt{Small}. Each residual polynomial
pair has finitely many polynomials as its matrix-entry
values, by local constancy and compactness. There are
only finitely many local flags and finitely many links.
All coefficients of these polynomials and the original
residual target therefore lie in one finite field
extension. Enlarge the coefficient ring to contain a
lift of that extension, and the finite coefficient
extension of the arithmetic seed and the ordinary links.
The inclusion of coefficient rings into the ambient
valuation ring is prescribed by \texttt{Amb}.

In characteristic zero, the complete DVR of the target
has finite residue field and fraction field finite over
$\mathbb Q_\ell$: its uniformizer has a positive integral
valuation on $\ell$, and the Cohen structure theorem
(or the usual complete-DVR degree calculation) gives
the finite extension. The prescribed embeddings agree
with the original coefficient map, and the compositum
is a finite extension of this complete valued field.
In characteristic $\ell$ the complete DVR is a power
series ring over its finite residue field. Extend its
constant field to the required residue field, with the
embedding into the ambient field given by the unique
finite-field lifts; all prescribed coefficient maps
factor through that finite field. In both cases the
enlarged DVR maps locally into the ambient valuation
ring and still realizes the target continuously.
We may further enlarge the coefficient ring and target
DVR to have exactly the same residue field as the
seed and all intermediate residual evaluations.
The characteristic-zero seed remains a lift of its
specified residue after this extension. Every object
now has the same base ring and residue field; apply
Theorem~\ref{thm:rh-four-space} to these actual lifts
and links.
\end{proof}
